% Appendix S5 — Implementation details, reproducibility, evaluation protocol.
% Auto-generated from env scan on 2026-04-24.
% Drop into sn-article.tex where an appendix section is expected.

\section*{Appendix S5 — Implementation details and evaluation protocol}
\label{appendix:s5}

\subsection*{S5.1 Software stack}

All experiments are reproduced on Python 3.12.10 with the following
library versions (pinned in \texttt{requirements.txt}):

\begin{itemize}\itemsep0pt
  \item \texttt{torch 2.11.0+cu128}, \texttt{torchvision 0.26.0+cu128},
        \texttt{torchaudio 2.11.0+cu128} (CUDA 12.8 runtime)
  \item \texttt{torch-geometric 2.7.0} — GraphSAGE embeddings
        (Section~\ref{sec:kg})
  \item \texttt{transformers 5.4.0}, \texttt{sentence-transformers 5.3.0}
        — Sentence-BERT \texttt{all-MiniLM-L6-v2} for content retrieval
  \item \texttt{faiss-cpu 1.13.2} — approximate nearest-neighbour index
  \item \texttt{implicit} — ALS collaborative filtering
  \item \texttt{xgboost 3.2.0}, \texttt{lightgbm 4.6.0} (used in
        experimental variants, not in the final pipeline)
  \item \texttt{scikit-learn 1.8.0}, \texttt{scipy 1.17.1},
        \texttt{numpy 1.26.4}, \texttt{pandas 3.0.1}
  \item \texttt{networkx 3.6.1} — prerequisite-graph manipulation
  \item \texttt{py-irt 0.1.1} — reference 3PL implementation used for
        calibration cross-checks
  \item \texttt{matplotlib 3.10.8} — figure rendering
\end{itemize}

\subsection*{S5.2 Hardware}

Single-GPU configuration: NVIDIA GeForce RTX 5050 Laptop GPU
(8\,151\,MiB VRAM, driver 573.13, CUDA 12.8). Host has 32\,GB system
RAM. No distributed training; every reported number is from a single
GPU run.

\subsection*{S5.3 Wall-clock per component}

Agent wall-clock per pipeline execution, averaged across seeds
$\{42, 123, 456, 789, 2024\}$ on XES3G5M with $n_{\text{students}}{=}6000$:

\begin{table}[h]
\centering
\small
\begin{tabular}{lrr}
\toprule
Component & Mean (s) & Std (s)\\
\midrule
Diagnostic (IRT calibration, EM) & 11.6   & 7.0  \\
Confidence (rule-based calibration) & 0.4   & 0.2  \\
Knowledge Graph + GraphSAGE + prerequisite mining & 347.9 & 104.0 \\
Prediction (transformer, 10.5 epochs mean) & 2{,}858.9 & 1{,}425.8 \\
Recommendation (index + priors) & 63.0  & 22.3 \\
Personalization (K-Means, 5 levels) & 3.0   & 1.1  \\
\bottomrule
\end{tabular}
\end{table}

Total end-to-end time per seed is approximately 55\,min, of which the
transformer training dominates (80--90\%). Evaluation on 900 held-out
test users adds $\sim$45\,s.

\subsection*{S5.4 Seeds and determinism}

Every experiment fixes a single global seed that is propagated to
(i) the user-level \texttt{train/val/test} split, (ii) PyTorch
initialisation (weights, dropout, DataLoader shuffle), (iii) NumPy
random state used by GraphSAGE link-prediction sampling,
(iv) Thompson-Sampling Beta priors, and (v) K-Means centroid
initialisation in the Personalization Agent. Propagation was audited
in Section~\ref{sec:repro} and is implemented by
\texttt{BaseAgent.global\_seed}; agent-level code may not silently
default to~42. The five seeds used throughout are
$\{42, 123, 456, 789, 2024\}$.

\subsection*{S5.5 Evaluation protocol}
\label{appendix:s5_eval}

For each test user, the chronologically-ordered interaction sequence
is split at the 30/70 boundary: the first 30\% of interactions form
the \emph{context} and seed the agent state (CAT estimate of
$\theta$, behavioural confidence history, prerequisite map); the
remaining 70\% form the \emph{evaluation horizon}. The Prediction
Agent emits a distribution over the 858 concepts; the ground-truth
top-$k$ label set for NDCG@10 / MRR / Precision@10 is the set of
concepts the user fails in the next 20 chronologically-ordered
interactions of the evaluation horizon (a re-attempted concept
counts once). Tag Coverage is the fraction of distinct concepts in
the union of all top-10 recommendations across test users,
normalised by the number of concepts that satisfy the macro-AUC
support criterion ($\geq 5$ positive and $\geq 5$ negative
interactions in the test horizon).

\subsection*{S5.6 Training objective}

The prediction head is optimised with a class-weighted focal binary
cross-entropy applied independently to each of the 858 concept
outputs. Let $\hat p_{i,t} = \sigma(z_{i,t})$ denote the predicted
failure probability for user-window $i$ on tag $t$, and let
$y_{i,t} \in \{0,1\}$ be the ground-truth label. Label smoothing is
applied to the targets:
\[
\tilde y_{i,t} = (1-\varepsilon)\, y_{i,t} + \tfrac{\varepsilon}{2},
\qquad \varepsilon = 0.05.
\]
The per-element focal BCE with positive-class weight $w_t$ is
\[
\ell_{i,t} = -\bigl[\,w_t\,\tilde y_{i,t} \log \hat p_{i,t}
                  + (1-\tilde y_{i,t}) \log (1-\hat p_{i,t})\,\bigr]
            \cdot (1-p^{\star}_{i,t})^{\gamma},
\]
where $p^{\star}_{i,t} = \hat p_{i,t}\,\tilde y_{i,t}
+ (1-\hat p_{i,t})(1-\tilde y_{i,t})$ is the probability assigned to
the (smoothed) correct class, $\gamma = 2.0$ is the focusing
parameter, and $w_t = \mathrm{clip}\!\bigl((1-r_t)/r_t,\,1,\,50\bigr)$
with $r_t$ the empirical positive rate of tag $t$ in the training
set. The total loss is $L = \tfrac{1}{NT}\sum_{i,t} \ell_{i,t}$.
Negative samples are not drawn explicitly: every concept that does
not appear in the next $\text{HORIZON}=20$ interactions is treated
as a negative for that window, which gives an average positive rate
of $\sim 1\%$ per tag and motivates the focal weighting.

\subsection*{S5.7 Leakage audit}

The 70/15/15 user-level split is fixed once per seed and shared by
all agents. IRT calibration (Section~\ref{subsec3c}), the GraphSAGE
embedding training, and prerequisite mining (Section~\ref{subsec3e})
operate exclusively on the 4{,}200 training users. Validation and
test users contribute neither to the IRT item parameters, nor to
the mastery sequences fed into the prerequisite mining heuristics,
nor to the GraphSAGE link-prediction objective. Tag embeddings
exposed to the Prediction Agent through warm initialisation
therefore contain no information about val/test responses, and the
reported AUC-ROC is not inflated by topology-mediated leakage.

\subsection*{S5.8 Reproduction command}

\begin{verbatim}
git clone https://github.com/RAli2505/Rec-system.git
cd Rec-system
pip install -r requirements.txt

# Full 5-seed main pipeline (~4-5 hours per seed on an RTX 5050):
python scripts/run_multi_seed.py --seeds 42 123 456 789 2024

# 5-seed inference-mode ablation (~75 min total):
python scripts/run_ablation_inference_5seeds.py

# Statistical tests + significance table:
python scripts/ablation_significance.py

# Confidence support + 6x6 matrix:
python scripts/confidence_support_analysis.py

# Per-tertile NDCG@10 for each config (seed 42, ~20 min):
python scripts/tertile_ablation_s42.py

# Post-hoc calibration (ECE / Brier / tertile):
python scripts/posthoc_calibration.py

# Sensitivity appendices (optional, ~6 h):
python scripts/sensitivity_kg_thresholds.py
python scripts/sensitivity_recommendation_weights.py
python scripts/sensitivity_context_window.py
\end{verbatim}
